\subsection*{横向比较：换用 Kafka 改变了哪些恢复位置}

前面的主方案已经闭环：购买库按状态轮询，RabbitMQ 按既定队列条件保留副本，持久核查接住暂时不能发布的工作。现在保持同一业务事件 EV1，只替换消息系统和读取方式，逐项比较哪些恢复依据随之改变。

若把 EV1 改写入 Kafka，还要判断应用恢复后的重发能否自动合并。这里先改变消息的保存模型：发布者向主题 \code{order-events} 追加记录，主题是一组有序分区的集合，每个分区按自己的日志位置保存记录。发布者可以用对象 O1 作为消息键，在固定分区规则与配置下，把该对象的记录选入同一分区。设 EV1 写入分区 2，记录取得该分区的位置，消息内容继续带本地事件号 EV1。日志位置标识这次写入，EV1 关联其业务事实~\cite{kafka-design}。

Kafka 也有名为幂等生产者的能力。其代理按生产者 ID 与发送序号识别协议重试，相关生产配置还限制确认、重试和在途请求数量~\cite{kafka-producer,kafka-design}。Outbox 进程重启以后重新调用一次 \code{send(EV1)}，会成为应用再次提交的记录。同一 EV1 写在消息内容中，代理就按普通消息内容保存，业务事件归并由相应处理协议完成。

换用 Kafka 时，发布器继续保存事件身份与本地恢复关系。消费方按组保存各分区的读取进度，日志则按主题配置保留，某个组推进位置后，其他组还可以读取同一条尚在保留中的记录。第十章从这份组进度推导故障后的重读边界。Kafka 事务可以覆盖它所管理的主题写入和相应消费进度，购买库 R 是另一提交域~\cite{kafka-design}。生产者协议去重减少协议重试造成的副本，应用在恢复时再次提交同一事件，则由稳定 EV1 关联这些记录。

\subsection*{日志订阅改变了哪一段工作}

轮询直接读取待发状态，另一种方式从数据库事务日志读取已提交事件，再转发给代理。事务日志保存数据库变更的提交依据，连接器持续读取并保存自己的进度。这减少了周期扫描的需要，同时把日志保留、连接器进度与重启恢复带入发布链。

Debezium Outbox Event Router 是把这类变更转换成事件消息的组件。其 \code{id} 用于事件身份，\code{aggregateid} 用作消息键，将同一业务对象路由到相应 Kafka 分区~\cite{debezium-outbox}。同一对象使用同一键，可以把它的消息放入同一分区。业务事务先形成不可变事件行，连接器按日志位置恢复读取，消费者再依据自己的处理进度提交效果。各段分别保存恢复位置，才把同一分区的顺序用于实际业务。

该组件对事件表有具体输入规则：预期业务事件以 \code{INSERT} 写入，\code{UPDATE} 按配置警告、报错或停止处理，\code{DELETE} 会被过滤~\cite{debezium-outbox}。本章轮询表会反复更新 \code{phase}、领取和受理状态。接入这个组件时，应把不可变事件行与轮询调度状态分开，让连接器读取前者的新增事件，并为连接器进度和日志保留建立恢复依据。

这两种变体分别改变代理的保存模型与发布器的事件读取路径。本章完整历史继续采用已经推导的轮询发布与 RabbitMQ：数据库保存待发状态，发布器领取事件，代理确认经过 PA 合流到 EV。
